% Section VII: Functorial TQFT and the Atiyah-Segal viewpoint
% Purpose: Formalize what the CS path integral is already doing
%
% Key message: "The functorial definition is not a separate subject pasted onto the physics;
%               it is the abstract codification of what the Chern-Simons path integral is already doing."
%
% Status: NEEDS TO BE WRITTEN (should be SHORT and clarifying)
% Sources:
% - Various references throughout existing files but not yet synthesized
% - This is mostly new writing
% - Keep to 2-3 pages maximum - introduce bordism category, state functor axioms, connect to CS
%
% Writing notes:
% - Introduce this AFTER reader has seen canonical quantization
% - Should feel inevitable, not arbitrary
% - Emphasize: this abstracts what we already computed in Section VI

\section{Functorial TQFT: Atiyah--Segal Axioms}

% TODO: Write new short section (~2-3 pages total)

\subsection{The Bordism Category}

% TODO: Define Bord_n briefly
% - Objects: closed (n-1)-manifolds
% - Morphisms: n-manifolds with boundary (bordisms)
% - Composition: gluing along boundaries
%
% Introduce bordism concept HERE (from topology_review.tex if content exists there)
% Keep it SHORT - we need definition, not full theory

\subsection{TQFT as a Functor}

% TODO: State Atiyah-Segal axioms
% Z: Bord_n → Vect_k is a symmetric monoidal functor
%
% Unpack what this means:
% - Objects Σ → vector spaces Z(Σ) = H(Σ)
% - Morphisms M: Σ₁ → Σ₂ → linear maps Z(M): H(Σ₁) → H(Σ₂)
% - Composition respects gluing: Z(M₁ ∪_Σ M₂) = Z(M₂) ∘ Z(M₁)
% - Symmetric monoidal: disjoint union ⊔ → tensor product ⊗

\subsection{Chern--Simons Theory as a TQFT}

% TODO: Connect back to CS theory explicitly
% "Everything we computed in Sections V-VI is an instance of this structure"
%
% - Surfaces → Hilbert spaces from quantization (Section VI)
% - 3-manifolds → path integral gives Z(M)
% - Gluing properties follow from path integral composition
%
% This is why CS theory is a "topological quantum field theory" -
% it realizes the functorial axioms, and correlation functions depend only on topology

% Keep this section SHORT and clarifying, not encyclopedic
% Purpose: legitimize the language, show it codifies what we already did
